\documentclass[12pt]{article}
\usepackage{amsmath, amssymb, amsthm}
\usepackage[utf8]{inputenc}
\usepackage[T1]{fontenc}
\usepackage{geometry}
\usepackage{hyperref}

\geometry{a4paper, margin=2.5cm}

\newtheorem{lemma}{Lemma}
\newtheorem{theorem}{Theorem}
\theoremstyle{remark}
\newtheorem*{remark}{Remark}

\title{\textbf{Bounds on the Number of Divisors}}
\author{0xHaru}
\date{\today}

% ---------------------------------------------------------------
\begin{document}
\maketitle

% ===============================================================
\section*{Preliminary Lemma}
% ===============================================================

\begin{lemma}
Let $d \in \mathbb{N}$ be a divisor of $n \in \mathbb{N}$.  Then at least one of $d$
or $k = n/d$ is less than or equal to $\sqrt{n}$.
\end{lemma}

\begin{proof}[Proof 1 (by contradiction)]
Since $d$ divides $n$, there exists an integer $k$ such that $d \cdot k = n$.

Suppose for contradiction that \emph{both} $d$ and $k$ are strictly greater
than $\sqrt{n}$:
\[
    d > \sqrt{n} \quad \text{and} \quad k > \sqrt{n}.
\]
Multiplying these two inequalities gives
\[
    d \cdot k > \sqrt{n} \cdot \sqrt{n} = n,
\]
which contradicts $d \cdot k = n$.  Therefore, at least one of $d$ or $k$
must satisfy $d \leq \sqrt{n}$ (or $k \leq \sqrt{n}$).
\end{proof}

\smallskip

\begin{proof}[Proof 2 (direct)]
Since $d$ divides $n$, there exists an integer $k = n/d$ such that
$d \cdot k = n$.  We consider two cases.

\medskip
\noindent\textit{Case $d \leq \sqrt{n}$.}\;
The thesis holds immediately, since $d$ itself is at most $\sqrt{n}$.

\medskip
\noindent\textit{Case $d > \sqrt{n}$.}\;
We show that $k < \sqrt{n}$:
\begin{alignat}{2}
    d &> \sqrt{n}
        &\qquad& \text{(assumption of this case)} \notag \\
    \frac{1}{d} &< \frac{1}{\sqrt{n}}
        &\qquad& \text{(taking reciprocals of positive quantities reverses the inequality)} \notag \\
    \frac{n}{d} &< \frac{n}{\sqrt{n}}
        &\qquad& \text{(multiplying both sides by $n > 0$)} \notag \\
    k = \frac{n}{d} &< \sqrt{n}
        &\qquad& \text{(simplifying $n/\sqrt{n} = \sqrt{n}$)} \notag
\end{alignat}

\medskip
\noindent In both cases at least one of $d$ or $k$ is less than or equal to
$\sqrt{n}$.
\end{proof}

\smallskip

\begin{remark}
This lemma tells us that every divisor of $n$ either is at most $\sqrt{n}$,
or its complementary divisor $n/d$ is at most $\sqrt{n}$.  In other words,
to enumerate \emph{all} divisors of $n$ it suffices to check only the
natural numbers up to $\sqrt{n}$.
\end{remark}

% ===============================================================
\section*{Main Theorem}
% ===============================================================

\begin{theorem}
For every $n \in \mathbb{N}$, the number of divisors $d(n)$ satisfies
\[
    d(n) \;\leq\; 2\sqrt{n}.
\]
\end{theorem}

\begin{proof}
Let $D$ denote the set of all positive divisors of $n$.  By definition
$|D| = d(n)$.

\smallskip
\noindent\textbf{Partition of $D$.}\;
We partition $D$ into three pairwise disjoint subsets:
\begin{align*}
    S_1 &= \bigl\{\, d \in D \;:\; d < \sqrt{n} \,\bigr\}, \\
    S_2 &= \bigl\{\, d \in D \;:\; d = \sqrt{n} \,\bigr\}, \\
    S_3 &= \bigl\{\, d \in D \;:\; d > \sqrt{n} \,\bigr\}.
\end{align*}
Since $S_1, S_2, S_3$ cover $D$ without overlap:
\[
    d(n) = |S_1| + |S_2| + |S_3|.
\]

% -----------------------------------------------------------
\bigskip
\noindent\textbf{Step 1: A bijection between $S_1$ and $S_3$.}

We show $|S_1| = |S_3|$ via the map
\[
    \varphi \;:\; S_1 \;\longrightarrow\; S_3, \qquad \varphi(d) = \frac{n}{d}.
\]
The key observation is that $\varphi$ is an \emph{involution}: applying it
twice returns the original element,
\[
    \varphi\bigl(\varphi(d)\bigr) = \varphi\!\left(\frac{n}{d}\right)
    = \frac{n}{n/d} = d.
\]
Because of this, $\varphi$ is its own inverse; to conclude it is a bijection
between $S_1$ and $S_3$ it suffices to check that it maps each set into the
other.

\medskip
\emph{$\varphi$ maps $S_1$ into $S_3$.}\;
Let $d \in S_1$, so $d$ divides $n$ and $d < \sqrt{n}$.
\begin{itemize}
    \item \emph{$n/d$ is a divisor of $n$:} Since $d \mid n$, we have
          $n/d \in \mathbb{N}$ and $d \cdot (n/d) = n$, so $n/d$ divides $n$.
    \item \emph{$n/d > \sqrt{n}$:} Since $d$ and $\sqrt{n}$ are positive,
          taking reciprocals of $d < \sqrt{n}$ reverses the inequality to
          $1/d > 1/\sqrt{n}$; multiplying both sides by $n > 0$ then gives
          $n/d > n/\sqrt{n} = \sqrt{n}$.
\end{itemize}
Hence $\varphi(d) = n/d \in S_3$.

\medskip
\emph{$\varphi$ maps $S_3$ into $S_1$.}\;
Let $d \in S_3$, so $d$ divides $n$ and $d > \sqrt{n}$.
\begin{itemize}
    \item \emph{$n/d$ is a divisor of $n$:} Since $d \mid n$, we have
          $n/d \in \mathbb{N}$ and $d \cdot (n/d) = n$, so $n/d$ divides $n$.
    \item \emph{$n/d < \sqrt{n}$:} Since $d$ and $\sqrt{n}$ are positive,
          taking reciprocals of $d > \sqrt{n}$ reverses the inequality to
          $1/d < 1/\sqrt{n}$; multiplying both sides by $n > 0$ then gives
          $n/d < n/\sqrt{n} = \sqrt{n}$.
\end{itemize}
Hence $\varphi(d) = n/d \in S_1$.

\medskip
Since $\varphi$ maps $S_1$ into $S_3$ and $S_3$ into $S_1$, and since
$\varphi \circ \varphi = \mathrm{id}$, the map $\varphi$ is a bijection
between $S_1$ and $S_3$.  Therefore $|S_1| = |S_3|$.

% -----------------------------------------------------------
\bigskip
\noindent\textbf{Step 2: Bounding $|S_1|$.}

Let $A(n) := |S_1|$.  Using $|S_1| = |S_3|$, we rewrite the divisor count as
\[
    d(n) = A(n) + |S_2| + A(n).
\]
By definition, $A(n)$ counts all naturals in $S_1$, each of which satisfies
$d < \sqrt{n}$.  The number of positive integers strictly less than $\sqrt{n}$
is at most $\lfloor\sqrt{n}\rfloor$, and $\lfloor\sqrt{n}\rfloor < \sqrt{n}$
always holds.\footnote{%
  If $n$ is not a perfect square then $\sqrt{n}$ is irrational, so
  $\lfloor\sqrt{n}\rfloor < \sqrt{n}$ directly.
  If $n$ is a perfect square, the positive integers strictly less than
  $\sqrt{n}$ number at most $\sqrt{n}-1 < \sqrt{n}$.}
Therefore
\[
    A(n) < \sqrt{n}.
\]

% -----------------------------------------------------------
\pagebreak
\noindent\textbf{Step 3: Final bounds.}

We distinguish two cases.

\medskip
\noindent\textit{Case A: $n$ is not a perfect square.}

$\sqrt{n} \notin \mathbb{N}$, so no divisor of $n$ can equal $\sqrt{n}$; hence
$S_2 = \varnothing$ and $|S_2| = 0$.  Therefore
\[
    d(n) = A(n) + 0 + A(n) = 2A(n) < 2\sqrt{n}.
\]

\medskip
\noindent\textit{Case B: $n$ is a perfect square.}

$\sqrt{n} \in \mathbb{N}$ and $\sqrt{n}$ divides $n$ (since
$\sqrt{n} \cdot \sqrt{n} = n$).  Thus $S_2$ contains exactly one element,
namely $\sqrt{n}$, so $|S_2| = 1$.  Therefore
\[
    d(n) = A(n) + 1 + A(n) = 2A(n) + 1 < 2\sqrt{n} + 1.
\]
Since $n$ is a perfect square both $d(n)$ and $2\sqrt{n}$ are integers, so
the strict inequality with $2\sqrt{n}+1$ is equivalent to
\[
    d(n) \leq 2\sqrt{n}.
\]

\medskip
\noindent In both cases we obtain $d(n) \leq 2\sqrt{n}$, which completes the
proof.
\end{proof}

\end{document}
